\section{Cross Section and Ratio Systematic Errors}
The systematic errors on cross sections and ratios of cross sections 
in this analysis are due to 
uncertainties on the number of selected background events, the incident 
neutrino flux, the number of targets in the detector, and the selection efficiencies. 
%$\epsilon$.
The sources of systematic uncertainties can be categorized into three groups:
%those from the neutrino 
beam flux prediction, 
neutrino and antineutrino interaction models 
and detector response.
The largest source of uncertainty is due to the beam flux.
\iffalse
Maybe write here that:  
The error estimations due to the beam flux and neutrino interaction modeling 
uncertainties utilized a reweighting approach, where one assign a weight 
for each event due (find a different word) to the property one studied
(flux, CCQE, etc) then one calculates the final observable 
and the estimated error is extract by comparison of the final observation 
with the nominal one.
\fi
\subsection{Beam flux uncertainty}
The beam flux uncertainty sources can be separated
into two categories: uncertainties 
of the hadronic interactions, in the graphite target and reinteractions in the
horn, and
T2K beamline inaccuracies. \\

The beam flux uncertainty is dominated by the
uncertainty on the modeling of the hadron interactions, including
uncertainties on the total proton-nucleus production cross
section, pion and kaon multiplicities, and secondary nucleon production. \\

The hadronic interactions in the target where the primary
proton beam first interacts and produces the majority of
the secondary pions is simulated by the FLUKA2011 package
which creates MC neutrino and antineutrino flux samples.
%[44,45]
%{
%[44] A. Ferrari, P. R. Sala, A. Fasso, and J. Ranft, Reports
%No. CERN-2005-010, No. SLAC-R-773, and No. INFNTC-
%05-11, 2005.
%[45] G. Battistoni, F. Cerutti, A. Fass�, A. Ferrari, S. Muraro, J.
%Ranft, S. Roesler, and P. R. Sala, AIP Conf. Proc. 896, 31
%(2007), we used FLUKA2008, which was the latest version
%at the time of this study. A new version, FLUKA2011, has
%been already released now and the comparison with data
%would be different.
%}
%Kinematic information for particles emitted from the target
%is transferred to the JNUBEAM simulation. 
%JNUBEAM is a MC simulation
%of the horn magnets, helium vessel, decay volume, beam dump, and muon monitor. 
%Hadronic interactions are modeled by the GCALOR model 
%[24,25]
%{
%[24] C. Zeitnitz and T. A. Gabriel, in Proceedings of
%International Conference on Calorimetry in High Energy
%Physics (Elsevier Science B.V., Tallahassee, FL, 1993).
%[25] A. Fasso, A. Ferrari, J. Ranft, and P. R. Sala, in Proceedings
%of the International Conference on Calorimetry in High
%Energy Physics, 1994.
%}
%in JNUBEAM. 
%The beam flux and energy spectra predications are tuned with an existing 
%hadron production data from NA61/SHINE measurements. 
%[30,31]
%{
%[30] N. Abgrall et al. (NA61/SHINE Collaboration), Phys. Rev.
%C 84, 034604 (2011).
%[31] N. Abgrall et al. (NA61/SHINE Collaboration), Phys. Rev.
%C 85, 035210 (2012).
%}
Uncertainties on the proton beam properties, horn current,
hadron production model and alignment are taken into
account to produce an energy-dependent systematic uncertainty
on the neutrino flux.
%These uncertainties are derived primarily from NA61/
%SHINE measurements 
%of pion and kaon production
%from proton interactions on a graphite target at the T2K 
%beam energy. 
%for areas of phase space uncovered by these
%measurements, additional external data sets are used. 
These uncertainties are propagated to the T2K neutrino beam flux
prediction by reweighting MC flux samples. 
The total proton-nucleus production cross section uncertainty
is adjusted to replicate discrepancies between NA61/SHINE measurements
%[30] 
and other external data sets. 
%[68-70]
%{
%[68] S. P. Denisov, S. V. Donskov, Yu. P. Gorin, R. N.
%Krasnokutsky, A. I. Petrukhin, Yu. D. Prokoshkin, and
%D. A. Stoyanova, Nucl. Phys. B61, 62 (1973).
%[69] G. Bellettini, G. Cocconi, A. N. Diddens, E. Lillethun, G.
%Matthiae, J. P. Scanlon, and A. M. Wetherell, Nucl. Phys.
%79, 609 (1966).
%[70] A. S. Carrol et al., Phys. Lett. 80B, 319 (1979).
%}
\\

%%.
%%The uncertainty
%%on secondary nucleon production is set by comparing the
%%predictions of the FLUKA2008 model with external
%%measurements [71,72]
%%{
%%[71] T. Eichten et al., Nucl. Phys. B44, 333 (1972).
%%[72] J. V. Allaby et al., CERN Technical Report No. 70-12, 1970.
%%}

%%Uncertainties in the operational conditions of the T2K beam-line 
%%are also considered. These include the proton beam
%%incident position on the target, the proton beam profile and
%%intensity, the angle between the beam and the ND280
%%detector, and the horn current, field, and alignment. \\
%%
%%The method of estimating these flux
%%uncertainties is described in Ref. [30]
%%{
%%[30] K. Abe et al. (T2K Collaboration), Phys. Rev. D 87, 012001
%%(2013).
%%}
%%.\\

%The beam flux uncertainties are given in a covariance matrix 
%for the FHC and RHC beam modes in terms of
%25 bins of neutrino energy from 0 to 30 GeV in 
%four neutrino flux modes
%of $\nu_{\mu}$,
%$\overline{\nu}_{\mu}$,
%$\nu_{e}$,
%and
%$\overline{\nu}_{e}$.
%The evaluation of the systematic cross section errors due to the flux uncertainties 
%was done both by analytic calculations
%and numerical toy experiments.

\iffalse
The effect of the flux uncertainties on individual cross sections
lead to $\approx$10\% errors whereas the ratio errors have ~3\% errors
due to correlated neutrino and antineutrino flux covariance errors. 
Table {\color{red} nnn} summarizes the systematic errors due to the beam flux 
uncertainties on the different results.
\fi

The flux smearing is done using toy MC data sets that are
based on the FHC and RHC beam flux uncertainty covariance matrices. 
The resulting $\pm1\sigma$ change in the cross section is taken as
the systematic error associated with the beam flux. 
These uncertainties on individual cross sections
lead to 9\% errors whereas the errors on the ratio are 4\%
due to correlated neutrino and antineutrino flux covariance errors. 
Table III summarizes the systematic errors due to the beam flux 
uncertainties on the cross sections and combinations of cross sections.
These results have been cross checked with analytic calculations.
The fractional errors on ratios have smaller errors due to
cancellations of correlated errors between the
neutrino and antineutrino modes.
\begin{table}
\caption{
\label{tab:BeamFluxSummary}
Summary table for one standard deviation errors due to beam flux uncertainties. (fractionl errors in $\%$)
}
\resizebox{\columnwidth}{!}{
\begin{ruledtabular}
\begin{tabular}{cccccc}
%\hline
%\hline
$\sigma(\overline{\nu})$ & $\sigma(\nu)$ & $R (\nu,\overline{\nu})$ 
& $A(\nu,\overline{\nu})$ & $\Sigma(\nu,\overline{\nu})$ & $\Delta (\nu,\overline{\nu})$ \\
\hline
 $\pm9.37$ & $\pm9.14$ & $\pm3.58$ & $\pm3.35$ & $\pm9.17$ & $\pm9.42$ \\
%\hline
%\hline
\end{tabular}
\end{ruledtabular}
}
\end{table}

%%Further details of
%%the neutrino beam flux prediction can be found in [27]
%%{
%%[27] K. Abe et al. (T2K Collaboration), Phys. Rev. D 87, 012001
%%(2013).
%%}
%%.

\subsection{Interaction model uncertainty}

The interaction model uncertainties were calculated by a data-driven method \citep{int-model} 
where the NEUT predictions were compared 
to external neutrino-nucleus data in the energy region relevant for T2K. 
Some of the NEUT model parameters are fitted and 
assigned mean and $1$ $\sigma$ error values that allow
for differences between NEUT 
and the external data. \\

%The NEUT parameters include axial mass values for
%quasi-elastic scattering and meson production via baryon resonances, 
%the Fermi momentum, the binding energy, 
%and a non-pion decay $\Delta$ parameter.
%The Spectral function parameter is introduced to take
%nto account the difference between the relativistic Fermi
%gas and Spectral function models. 
%The Fermi gas model is in the standard NEUT simulation.
%The more sophisticated Spectral function model 
%uses electron scattering data 
%and is expected to be a better representation of nuclear interactions.
%In addition, uncertainties on the final state
%interactions of the pions and nucleons within the nuclear
%medium are included.


\begin{table*}
\caption{
\label{tab:PhysicsModelPiFSISummary}
Summary table for physics model uncertainties for restricted
phase space measurements (fractional errors in $\%$).
}

%\resizebox{\columnwidth}{!}{
%\begin{ruledtabular}
\begin{tabular}{lcccccc}
\hline
\hline
Parameter & 
$\sigma(\overline{\nu})$ & $\sigma(\nu)$ & $R (\nu,\overline{\nu})$ & $A(\nu,\overline{\nu})$ & $\Sigma(\nu,\overline{\nu})$ & $\Delta (\nu,\overline{\nu})$ \\
%$\sigma(\overline{\nu})$ & $\sigma(\nu)$ & $\sigma(\overline{\nu})/\sigma(\nu)$ & $(\sigma(\nu)-\sigma(\overline{\nu}))/(\sigma(\nu)+\sigma(\overline{\nu}))$ & $\sigma(\nu)+\sigma(\overline{\nu})$ & $\sigma(\nu)-\sigma(\overline{\nu})$ \\
\hline
$M^{QE}_A$ 
& $\pm0.51$ & $\pm0.14$ & $\pm0.37$ & $\pm0.32$ & $\pm0.24$ & $\pm0.08$ \\
$p_F\ (^{12}$C) 
& $\pm0.01$ & $\pm0.02$ & $\pm0.01$ & $\pm0.01$ & $\pm0.02$ & $\pm0.02$ \\
$p_F\ (^{16}$O) 
& $0$ & $\pm0.01$ & $0$ & $0$ & $\pm0.01$ & $\pm0.01$ \\
MEC norm\ $(^{12}$C) 
& $\pm0.30$ & $\pm0.44$ & $\pm0.14$ & $\pm0.12$ & $\pm0.40$ & $\pm0.52$ \\
MEC norm\ $(^{16}$O) 
& $\pm0.18$ & $\pm0.24$ & $\pm0.06$ & $\pm0.05$ & $\pm0.22$ & $\pm0.27$ \\
$E_B\ (^{12}$C) 
& $\pm0.01$ & $\pm0.01$ & $\pm0.02$ & $\pm0.02$ & $0$ & $\pm0.02$ \\
$E_B\ (^{16}$O) 
& $\pm0.01$ & $\pm0.01$ & $\pm0.02$ & $\pm0.02$ & $0$ & $\pm0.02$ \\
$C_5^A$(0) 
& $\pm0.70$ & $\pm0.46$ & $\pm0.24$ & $\pm0.21$ & $\pm0.53$ & $\pm0.32$ \\
$M_A^{1\pi}$ 
& $\pm0.99$ & $\pm0.28$ & $\pm0.75$ & $\pm0.65$ & $\pm0.44$ & $\pm0.21$ \\
$I=\frac{1}{2}$ Bkg 
& $\pm0.29$ & $\pm0.21$ & $\pm0.08$ & $\pm0.07$ & $\pm0.23$ & $\pm0.17$ \\
$\nu_e/\nu_{\mu}$
& $\pm0.02$ & $0$ & $\pm0.01$ & $\pm0.01$ & $\pm0.01$ & $0$   \\
CC Other shape 
& $\pm0.65$ & $\pm0.70$ & $\pm0.06$ & $\pm0.79$ & $\pm0.06$ & $\pm0.75$ \\
CC Coherent 
& $\pm0.01$ & $\pm0.01$ & $0$ & $\pm0.05$ & $\pm0.69$ & $\pm0.73$ \\
NC Coherent 
& $0$ & $0$ & $0$ & $0$ & $0$ & $0$ \\
NC Other 
& $\pm1.28$ & $\pm0.39$ & $\pm0.89$ & $\pm0.77$ & $\pm0.63$ & $\pm0.14$ \\
$\pi$ FSI
& $\pm0.16$ & $\pm0.19$ & $\pm0.11$ & $\pm0.09$ & $\pm0.18$ & $\pm0.23$ \\
MEC norm Other & $\pm0.08$ & $\pm0.15$ & $\pm0.07$ & $\pm0.20$ & $\pm0.13$ & $\pm0.20$ \\
\hline
Total & $\pm2.13$ & $\pm1.16$ & $\pm1.56$ & $\pm1.36$ & $\pm1.31$ & $\pm1.32$\\
\hline
\hline
\end{tabular}
%\end{ruledtabular}
%}
\end{table*}


\begin{table*}
\caption{
\label{tab:DetectorResponseSummary} 
Summary table for detector response uncertainties (fractional errors in $\%$).
}

%\resizebox{\columnwidth}{!}{
%\begin{ruledtabular}
\begin{tabular}{lcccccc}
\hline
\hline
Parameter &
$\sigma(\overline{\nu})$ & $\sigma(\nu)$ & $R (\nu,\overline{\nu})$ & $A(\nu,\overline{\nu})$ & $\Sigma(\nu,\overline{\nu})$ & $\Delta (\nu,\overline{\nu})$ \\
%& $\sigma(\overline{\nu})$ & $\sigma(\nu)$ & $\sigma(\overline{\nu})/\sigma(\nu)$ & $(\sigma(\nu)-\sigma(\overline{\nu}))/(\sigma(\nu)+\sigma(\overline{\nu}))$ & $\sigma(\nu)+\sigma(\overline{\nu})$ & $\sigma(\nu)-\sigma(\overline{\nu})$ \\
\hline
TPC tracking Efficiency
& $\pm0.37$ & $\pm0.32$ & $\pm0.04$ & $\pm0.04$ & $\pm0.34$ & $\pm0.29$ \\
Charge misidentification
& $\pm0.37$ & $\pm0.32$ & $\pm0.04$ & $\pm0.04$ & $\pm0.34$ & $\pm0.29$ \\
Sand/Rock muon interference
& $\pm1.45$ & $\pm2.20$ & $\pm0.74$ & $\pm0.70$ & $\pm1.99$ & $\pm2.70$ \\
Fiducial mass
& $\pm0.96$ & $\pm0.96$ & $0$ & $0$ & $\pm1.36$ & $\pm1.36$ \\
Fiducial volume boundaries
& $\pm0.13$ & $\pm0.97$ & $\pm0.83$ & $\pm1.39$ & $\pm0.77$ & $\pm0.74$ \\
%Highest momentum track in bunch %GUESS this is OK???
%& $\pm0.13$ & $0$ & $\pm0.13$ & $\pm0.11$ & $\pm0.04$ & $\pm0.08$ \\
\hline
Total & $\pm1.82$ & $\pm2.63$ & $\pm1.11$ & $\pm1.02$ & $\pm2.58$ & $\pm3.35$\\
\hline
\hline
\end{tabular}
%\end{ruledtabular}
%}
\end{table*}


The CCQE model in NEUT is based on the Llewellyn-Smith neutrino-nucleon scattering model \citep{LSmith}
% C.~H.~Llewellyn Smith,``Neutrino Reactions at Accelerator Energies’’, Phys.\ Rept.\  {\bf 3}, 261 (1972).
with a dipole axial form factor and the BBBA05 vector form factors \citep{RBradford}.
% R.~Bradford, A.~Bodek, H.~S.~Budd and J.~Arrington,``A New parameterization of the nucleon elastic form-factors,'' Nucl.\ Phys.\ Proc.\ Suppl.\  {\bf 159}, 127 (2006)
The NEUT generator uses the Smith-Moniz RFG model \citep{RSmith}
% R.~A.~Smith and E.~J.~Moniz, ``Neutrino Reactions On Nuclear Targets,'' Nucl.\ Phys.\ B {\bf 43}, 605 (1972); Erratum: [Nucl.\ Phys.\ B {\bf 101}, 547 (1975)]. 
and includes an implementation of both the random phase approximation (RPA) correction \citep{JNieves}
% J.~Nieves, I.~Ruiz Simo and M.~J.~Vicente Vacas, ``Inclusive Charged--Current Neutrino--Nucleus Reactions,'' Phys.\ Rev.\ C {\bf 83}, 045501 (2011)
and the 2p2h Nieves model \citep{JNieves}.
% J.~Nieves, I.~Ruiz Simo and M.~J.~Vicente Vacas, ``Inclusive Charged--Current Neutrino--Nucleus Reactions,'' Phys.\ Rev.\ C {\bf 83}, 045501 (2011).
% R.~Gran, J.~Nieves, F.~Sanchez and M.~J.~Vicente Vacas, ``Neutrino-nucleus quasi-elastic and 2p2h interactions up to 10 GeV,'' Phys.\ Rev.\ D {\bf 88}, no. 11, 113007 (2013)
The NEUT resonant pion production is based on the Rein-Sehgal model \citep{DRein}
% D.~Rein and L.~M.~Sehgal, ``Neutrino Excitation of Baryon Resonances and Single Pion Production,'' Annals Phys.\  {\bf 133}, 79 (1981).
with updated form factors from Ref. \citep{KGraczyk}.
% K.~M.~Graczyk and J.~T.~Sobczyk, ``Form Factors in the Quark Resonance Model,'' Phys.\ Rev.\ D {\bf 77}, 053001 (2008); Erratum: [Phys.\ Rev.\ D {\bf 79}, 079903 (2009)]
The DIS model
%[ {\color{red} (missing reference)n } ]
used in NEUT includes both the structure function from Ref. \citep{MGluck}
% M.~Glück, E.~Reya and A.~Vogt, ``Dynamical parton distributions revisited,'' Eur.\ Phys.\ J.\ C {\bf 5}, 461 (1998)
and the Bodek-Yang correction \citep{ABodek}.
% A.~Bodek and U.~K.~Yang, ``Modeling neutrino and electron scattering cross-sections in the few GeV region with effective LO PDFs,'' AIP Conf.\ Proc.\  {\bf 670}, 110 (2003) 
The NEUT MC generator includes various model parameters to describe the different models, uncertainties and approximations. 
The axial mass $M_A^{QE}$ 
was set to 1.21 GeV/c$^2$ based on the Super-Kamiokande atmospheric data and the K2K data.
%The NEUT MC generator uses the
%Llewellyn-Smith CCQE neutrino-nucleon cross-section
%with the relativistic Fermi gas model by Smith and Moniz. 
%From tuning to Super-Kamiokande atmospheric data and
%K2K data, the nominal axial mass $M^{QE}_A$ was set to 1.21 GeV/c$^2$. 
The 1$\sigma$ error on $M^{QE}_A$ was set to 0.41 GeV/c$^2$.
The large uncertainty on this parameter is due to the disagreements between
recent experimental measurements and bubble chamber results\citep{maqe}.
The Fermi gas momentum parameter ($p_F$) values and their errors are set to 
223 MeV/c and 225 MeV/c for Carbon and Oxygen respectively 
with both errors set to $\pm$12.7 MeV/c. 
The Fermi gas binding energy ($E_B$) parameter was set to 
25 MeV and 27 MeV for Carbon and Oxygen respectively 
with both errors set to $\pm$9 MeV. 
%The 2p2h model from Nieves $et\ al.$ \citep{JNieves-2,nieves-1} 
%[34, 35] 
%{
%[34] J. Nieves, I. Ruiz Simo, and M. J. Vicente Vacas. Inclusive charged-current
%1077 neutrino-nucleus reactions. Phys. Rev. C, 83:045501, Apr 2011.
%[35] R. Gran, J. Nieves, F. Sanchez, and M. J. Vicente Vacas. Neutrino-nucleus
%1455 quasi-elastic and 2p2h interactions up to 10 GeV. 2013.
%},ni
%has been implemented in NEUT.
%meson exchange current (MEC) 
The Nieves model 2p2h normalization to  $1\pm1$ 
%100\%+/-100\% 
for both Carbon and Oxygen,
the resonant pion production model in NEUT used the Graczyk and Sobczyk 
%[44]
%{
%Krzysztof M. Graczyk and Jan T. Sobczyk. Form Factors in the Quark Resonance Model. Phys. Rev., D77:053001, 2008.
%}
form factors $C^A_5(0)$ 
and the $I=\frac{1}{2}$ background scale 
were set to $1.01\pm0.12$ and $1.20\pm0.20$ respectively. 
The nominal axial mass $M^{RES}_A$ was set to $0.95\pm0.15 {\rm\  GeV/c}^2$.
Additional uncertainties are $\nu_e/\nu_\mu$ cross section factor 
that was set to $1.00\pm0.02$. Both CC and NC coherent uncertainties 
based on the Rein-Sehgal model
were set to $1\pm1$ and $1.0\pm0.3$ respectively. 
Moreover, for CC and NC interactions, additional scale factors
were set to $0.0\pm0.4$ and $1.0\pm0.3$ respectively.
In addition the CC other is an energy dependent factor\citep{tpc-analysis-paper}
and the NC other is a normalization factor.
The $\pi$ Final State Interaction (FSI) uncertainties 
are tuned to a pion-nucleus scattering data, and 
other smaller corrections were included  \citep{int-model}.

Variation of model parameters within their errors ($\pm1\sigma$) 
was used to estimate their effect on the final observables in order to
determine final measurement uncertainties.
A summary of the parameters and their effects on the overall normalization 
are shown in Table \ref{tab:PhysicsModelPiFSISummary}.

\subsection{Detector response uncertainty}
The detector response uncertainty studies 
%performed in the analysis were driven by
used data samples supported with MC samples
and measurements of the target weight.
The three dominant detector response 
systematic uncertainties %originate from the
are caused by the fiducial volume boundaries, 
the sand/rock muon interactions 
and the mass of the target in the fiducial volume.
There were small uncertainties from reconstruction and
charge misidentification from the TPC measurements. 
All the sources of detector response errors considered 
in the analysis are given 
in Table \ref{tab:DetectorResponseSummary}.

The fiducial volume systematics were
estimated by varying its  boundaries.
%the dependence of the individual cross sections 
%and the other final measurements on the fiducial volume definition .
The sand/rock muon interactions occurring upstream and 
in the surrounding ND280 volume
could create tracks passing through the P$\emptyset$D and TPC detectors,
mimicking a CCINC event. 
Another source of detector systematics was the mass of the target 
in the fiducial volume. The uncertainty due to the fiducial mass 
was conservatively estimated to be 0.96\%
from the measured mass of the detector material during construction and
the water mass measured during filling the water bags.


















